\section{Kết quả thực nghiệm và đánh giá}

\subsection{Kết quả xây dựng hệ thống}

\subsubsection{Data Warehouse}

Hệ thống Kho dữ liệu Global Superstore đã được triển khai và vận hành ổn định trên nền tảng PostgreSQL 17 đóng gói trong Docker container, khởi tạo đầy đủ 4 lược đồ chuyên biệt (\texttt{bronze}, \texttt{silver}, \texttt{gold}, \texttt{dss}) với tổng cộng 15 bảng vật lý và 3 bảng tổng hợp Materialized Views thông qua lệnh điều phối duy nhất \texttt{python -m src.pipeline.cli --init-db}. Tại tầng Bronze, hệ thống đã tiếp nhận toàn vẹn 51.290 bản ghi giao dịch gốc vào bảng \texttt{bronze.raw\_orders}, trích xuất chuẩn xác 632 bản ghi mã bưu chính vào \texttt{bronze.postal\_code} và 149 ánh xạ quốc gia -- thị trường vào \texttt{bronze.market}, đồng thời lưu vết kiểm toán SHA-256 vào \texttt{bronze.ingestion}. Tại tầng Silver, bảng \texttt{silver.orders} đã chuẩn hóa thành công 100\% bản ghi sang kiểu dữ liệu định lượng và ngày tháng chuẩn ISO mà không phát sinh bất kỳ lỗi vi phạm ràng buộc miền giá trị nào. Tại tầng Gold, mô hình hình sao hoàn chỉnh đã được thiết lập với bảng sự kiện trung tâm \texttt{gold.fact\_sales} (51.290 dòng) kết nối chặt chẽ qua các khóa thay thế số nguyên với 5 bảng chiều gồm \texttt{gold.dim\_date} (2.922 ngày lịch), \texttt{gold.dim\_customer} (1.590 khách hàng), \texttt{gold.dim\_product} (10.292 sản phẩm), \texttt{gold.dim\_geography} (3.847 vị trí địa lý) và \texttt{gold.dim\_ship\_mode} (4 phương thức giao vận), cùng 3 Data Marts tính toán sẵn (\texttt{mv\_sales\_mart}, \texttt{mv\_product\_mart}, \texttt{mv\_customer\_mart}).

\subsubsection{Dashboard và giao diện truy vấn}

Hệ thống trình diễn và truy vấn dữ liệu được hiện thực hóa song song trên hai nền tảng nhằm phục vụ đa dạng nhu cầu của người sử dụng. Trên nền tảng ứng dụng Web, giao diện điều hành tương tác được xây dựng bằng ngôn ngữ Python và thư viện Streamlit (\texttt{src/app/main.py}) kết hợp cùng dịch vụ RESTful API trên nền tảng FastAPI (\texttt{src/app/api/main.py}). Giao diện Web được chia thành bốn phân trang trực quan kết nối trực tiếp với cơ sở dữ liệu PostgreSQL: Phân trang tổng quan hiển thị 5 thẻ chỉ số KPI toàn cục, hệ thống cảnh báo điểm nóng rủi ro và biểu đồ kết hợp cột -- đường trục kép phản ánh diễn biến Doanh thu, Lợi nhuận và Số lượng đơn hàng qua 48 tháng liên tục; các phân trang còn lại cung cấp công cụ tương tác thời gian thực với ba mô-đun của hệ hỗ trợ quyết định.

Song song với ứng dụng Web, bộ tài nguyên báo cáo quản trị trên phần mềm Power BI Desktop đã được xây dựng hoàn chỉnh trong thư mục \texttt{reports/powerbi/}, bao gồm tệp cấu hình giao diện màu sắc chuẩn nhận diện (\texttt{hcmut\_theme.json}), tập lệnh kết nối trực tiếp cơ sở dữ liệu Power Query M (\texttt{02\_POWER\_QUERY\_M\_SCRIPTS.m}), thư viện hơn 20 công thức đo lường động DAX (\texttt{03\_DAX\_MEASURES.dax}) như \texttt{Total Sales}, \texttt{Total Profit}, \texttt{Profit Margin \%}, \texttt{Sales YoY Growth \%}, \texttt{Pareto \%}, \texttt{Negative Profit Orders Count} và tiện ích xuất dữ liệu sạch (\texttt{src/app/export\_for\_powerbi.py}). Báo cáo được tổ chức thành ba trang phân tích chuyên sâu: Trang 1 (\textit{Executive Sales Performance}) theo dõi hiệu suất tài chính theo thời gian và bản đồ địa lý toàn cầu; Trang 2 (\textit{Product Portfolio \& Pareto ABC}) phân tích ma trận định vị sản phẩm và quy luật Pareto 80/20; Trang 3 (\textit{Customer Analytics \& Cross-Selling}) phân tích cơ cấu phân khúc khách hàng và bảng đề xuất bán chéo sản phẩm.

\subsubsection{Hệ thống hỗ trợ quyết định (DSS)}

Phân hệ DSS đã được hiện thực hóa đầy đủ với ba mô-đun thuật toán hoạt động độc lập qua giao diện dòng lệnh (CLI) hoặc gọi trực tiếp từ giao diện Web và API. Mô-đun dự báo chuỗi thời gian (\texttt{src/dss/predictive/run.py}) đã huấn luyện thành công mô hình San bằng mũ có xu thế trên chuỗi 48 tháng lịch sử, lưu trữ kết quả đánh giá vào bảng \texttt{dss.forecast\_run} và xuất 6 điểm dự báo doanh thu nửa đầu năm 2015 kèm cận dưới/cận trên của khoảng tin cậy 95\% vào bảng \texttt{dss.forecast\_result}. Mô-đun khai phá giỏ hàng (\texttt{src/dss/association/run.py}) đã phân tích 25.035 đơn đặt hàng và lưu trữ bền vững vào bảng \texttt{dss.association\_rule} cả hai cấp độ tri thức: 14 luật kết hợp ở cấp độ tiểu ngành hàng và 27 luật kết hợp ở cấp độ sản phẩm SKU chi tiết. Mô-đun mô phỏng kịch bản What-If (\texttt{src/dss/what\_if/simulate.py}) đã tích hợp hoàn chỉnh khả năng truy vấn động số liệu gốc từ kho dữ liệu Gold theo từng bộ lọc thị trường -- ngành hàng, giải hệ phương trình co giãn cầu theo giá trong thời gian thực và đưa ra thông điệp khuyến nghị hành động ba mức độ rõ ràng.

\subsection{Đánh giá Data Warehouse và ETL}

Để đánh giá toàn diện chất lượng của Kho dữ liệu và quy trình ETL, nhóm thực hiện xem xét trên bốn tiêu chí kỹ thuật cốt lõi gồm tính chính xác dữ liệu, tính lũy đẳng và khả năng nạp gia tăng, hiệu năng xử lý đường ống và tốc độ phản hồi truy vấn OLAP.

Về tính chính xác và toàn vẹn dữ liệu, kết quả thực thi bộ kiểm định tự động tại cả hai tầng Silver và Gold đạt tỷ lệ vượt qua 100\% trên toàn bộ các bài kiểm tra. Không có bất kỳ bản ghi nào bị rỗng khóa bắt buộc, không có bản ghi trùng lặp \texttt{row\_id} và không tồn tại khóa ngoại mồ côi giữa bảng sự kiện \texttt{gold.fact\_sales} với năm bảng chiều. Phép đối soát tổng hợp tài chính giữa tầng Silver và tầng Gold khớp nhau tuyệt đối trên cả 51.290 bản ghi với tổng doanh thu đạt $12.642.501,91\text{ USD}$, tổng lợi nhuận đạt $1.467.457,29\text{ USD}$ và tổng số lượng hàng hóa đạt $178.312\text{ đơn vị}$. Các ca biên phức tạp của tập dữ liệu gốc như ngày tháng đa định dạng, 80,51\% mã bưu chính khuyết thiếu và hiện tượng thành phố thuộc nhiều khu vực thị trường đều đã được xử lý triệt để mà không phải loại bỏ bất kỳ giao dịch hợp lệ nào.

Về tính lũy đẳng (Idempotency) và khả năng nạp gia tăng (Incremental Loading), hệ thống chứng tỏ độ tin cậy cao trong vận hành thực tế. Nhờ cơ chế đối chiếu mã băm SHA-256 trong bảng \texttt{bronze.ingestion} kết hợp mệnh đề \texttt{ON CONFLICT DO NOTHING / DO UPDATE} tại tầng Silver và Gold, việc kích hoạt lại đường ống ETL nhiều lần trên cùng một tệp dữ liệu hoàn toàn không làm phát sinh bản ghi dư thừa. Khi kiểm thử nạp gia tăng với lô 500 bản ghi mới của Quý 1/2015, cơ chế cờ đánh dấu \texttt{transform = FALSE} giúp hệ thống chỉ xử lý đúng 500 bản ghi mới phát sinh trong thời gian chưa đầy 3 giây, nâng tổng số dòng trong \texttt{gold.fact\_sales} lên 51.790 dòng một cách chính xác.

Về hiệu năng truy vấn phân tích OLAP, việc chuyển đổi từ bảng phẳng sang mô hình hình sao sử dụng khóa thay thế số nguyên kết hợp đánh chỉ mục B-Tree trên toàn bộ các cột khóa ngoại đã mang lại tốc độ truy vấn vượt trội. Toàn bộ 7 truy vấn OLAP phức tạp (bao gồm các phép toán nặng như \texttt{GROUP BY ROLLUP} trên 4 chiều, \texttt{GROUP BY CUBE} sinh 8 tổ hợp chiều và các hàm cửa sổ \texttt{DENSE\_RANK}, \texttt{LAG}) khi chạy trực tiếp trên 51.290 dòng của bảng sự kiện đều hoàn tất trong khoảng thời gian từ $16,1\text{ ms}$ đến $49,2\text{ ms}$ (trung bình dưới $30\text{ ms}$). Đặc biệt, khi truy vấn thông qua 3 bảng Materialized Views đã được tổng hợp trước, thời gian phản hồi giảm xuống dưới $5\text{ ms}$, đảm bảo trải nghiệm mượt mà tức thì khi người dùng thao tác lọc chéo trên giao diện Web và Power BI.

\subsection{Đánh giá DSS}

Tính hiệu quả của Hệ hỗ trợ quyết định được đánh giá dựa trên độ chính xác của các mô hình định lượng và giá trị thực tiễn của các tri thức được khám phá đối với hoạt động quản trị tại Global Superstore. Đối với mô-đun dự báo chuỗi thời gian (Predictive Analytics), việc phân tách tập huấn luyện (2011--2013) và tập kiểm thử (2014) theo đúng trình tự thời gian giúp đánh giá trung thực năng lực dự báo. Mặc dù doanh số năm 2014 có biên độ dao động mùa vụ rất lớn giữa các tháng đầu năm và cuối năm, mô hình San bằng mũ kết hợp xu thế tuyến tính vẫn bám sát quỹ đạo tăng trưởng chung và đạt sai số phần trăm tuyệt đối trung bình $\text{MAPE} = 14,72\%$ ($\text{MAE} = 48.591,40\text{ USD}$). So với mô hình trung bình trượt cơ sở vốn bị trễ pha khi chuỗi thời gian có xu thế dốc lên, mô hình được chọn phản ánh chính xác tốc độ tăng trưởng trung bình hàng tháng và cung cấp dải tin cậy 95\% hợp lý cho 6 tháng đầu năm 2015.

Đối với mô-đun khai phá luật kết hợp (Association Rule Mining), việc thiết kế chiến lược khai phá hai cấp độ đã giải quyết triệt để bài toán thưa dữ liệu của danh mục 10.292 sản phẩm. Ở cấp độ tiểu ngành hàng (\texttt{sub\_category}), 14 luật tìm được đều có độ hỗ trợ từ 0,58\% đến 1,56\% và độ tin cậy lên tới 23,10\%, cung cấp cơ sở vững chắc để bố trí các gian hàng liên quan trên trang thương mại điện tử (như gợi ý Nhãn dán, Phong bì và Tủ sách ngay khi khách hàng xem danh mục Máy sao chụp). Ở cấp độ sản phẩm chi tiết (\texttt{product}), việc sử dụng ngưỡng đếm tuyệt đối $\text{min\_count} = 3$ đã giúp phát hiện 27 cặp sản phẩm có hệ số Lift rất cao (từ $31,91$ đến $315,57$), phục vụ trực tiếp cho tính năng đề xuất mua kèm tự động tại bước kiểm tra giỏ hàng.

Đối với mô-đun mô phỏng kịch bản What-If, giá trị lớn nhất mà mô-đun mang lại là việc kết nối trực tiếp lý thuyết độ co giãn cầu theo giá với số liệu biên lợi nhuận thực tế của từng thị trường trong kho dữ liệu. Thực nghiệm trên ngành hàng Nội thất và tiểu ngành hàng Bàn làm việc (\texttt{Tables}) đã chứng minh định lượng vì sao việc tăng chiết khấu tại các nhóm hàng có biên lợi nhuận gốc mỏng ($< 10\%$) lại gây ra sự sụt giảm lợi nhuận ròng nghiêm trọng dù doanh số bán hàng vẫn tăng, từ đó giúp ban giám đốc tránh được các chiến dịch khuyến mãi gây lỗ vốn.

\subsection{Hạn chế và hướng phát triển}

Bên cạnh những kết quả đã đạt được, hệ thống hiện tại vẫn tồn tại một số hạn chế nhất định xuất phát từ đặc thù của tập dữ liệu nguồn và phạm vi triển khai trong khuôn khổ môn học. Thứ nhất, dữ liệu lịch sử của Global Superstore chỉ bao gồm các giao dịch bán hàng đã hoàn tất mà không chứa dữ liệu về số lượng hàng tồn kho thực tế tại từng kho bãi (Inventory on hand), giá vốn nhập hàng chi tiết theo từng nhà cung cấp hoặc chi phí tiếp thị quảng cáo (Marketing Spend), do đó hệ thống chưa thể mở rộng sang các bài toán tối ưu hóa điểm đặt hàng lại (Reorder Point -- ROP) hay tính toán chi phí thu hút khách hàng (CAC). Thứ hai, ở mô-đun dự báo chuỗi thời gian, công thức dự báo điểm ở các bước xa ($h > 1$) hiện đang ngoại suy theo độ dốc xu thế tuyến tính kết hợp làm mượt mũ đơn giản nên chưa mô hình hóa đầy đủ thành phần dao động hình sin của chu kỳ mùa vụ 12 tháng vào các điểm dự báo xa. Thứ ba, quy trình nạp dữ liệu hiện tại vẫn vận hành theo phương thức xử lý theo lô (Batch Processing) kích hoạt qua dòng lệnh thay vì tự động lắng nghe sự kiện theo luồng thời gian thực.

Từ những hạn chế trên, nhóm đề xuất bốn hướng phát triển cụ thể cho hệ thống trong tương lai. Hướng phát triển thứ nhất là tích hợp công cụ điều phối quy trình công việc tự động như Apache Airflow hoặc Prefect để lập lịch chạy pipeline nạp gia tăng định kỳ hàng đêm kèm theo cơ chế gửi cảnh báo tự động khi phát hiện dữ liệu bất thường. Hướng phát triển thứ hai là nâng cấp mô-đun dự báo chuỗi thời gian bằng cách tích hợp các mô hình có thành phần mùa vụ tường minh như SARIMA (Seasonal ARIMA), Holt-Winters ba thành phần (Triple Exponential Smoothing) hoặc mô hình học máy Gradient Boosting (XGBoost, LightGBM) với các đặc trưng trễ thời gian (Lag features) để giảm sai số MAPE xuống dưới 8\%. Hướng phát triển thứ ba là mở rộng chiến lược Chiều biến đổi chậm tại bảng \texttt{gold.dim\_customer} và \texttt{gold.dim\_product} từ SCD Type 1 sang SCD Type 2 hoàn chỉnh với các cột \texttt{effective\_from}, \texttt{effective\_to} và \texttt{is\_current} nhằm theo dõi lịch sử chuyển đổi phân khúc của khách hàng qua nhiều năm. Hướng phát triển thứ tư là kết hợp kết quả phân khúc khách hàng RFM tại \texttt{gold.mv\_customer\_mart} với các luật kết hợp giỏ hàng để xây dựng hệ thống gợi ý cá nhân hóa (Personalized Recommendation System) cho từng nhóm khách hàng mục tiêu.
